\chapter{Descente icosaédrique \texorpdfstring{\(A_5/C_5\)}{A₅/C₅} et sélection du support gravitationnel à cinq composantes}
\label{chap:parte-iv-63}

La droite gravitationnelle reçoit une échelle chronologique puis descend au quotient icosaédrique. Le projecteur irréductible de rang \(5\) sélectionne le support tensoriel sans encore supposer de réalisation massive ou radiative.

\section{Constante gravitationnelle intrinsèque et représentation tensorielle}
\label{rm:ch:gravity}

\subsection{Chaîne constitutive du secteur gravitationnel : échelle chronologique \texorpdfstring{\(108\)}{108}, réduction tensorielle, holonomie et évaluation métrologique de \texorpdfstring{\(G\)}{G}}

La gravité s'organise ici en quatre strates qu'il ne faut pas fusionner :

\begin{equation}
\label{rm:eq:gravity-causal-chain}
\begin{aligned}
\text{structure chronologique CT108 et droites d'action}
&\longrightarrow \text{sélecteur de périodicité }\theta_{108}^{\rm circ}
\longrightarrow G_{108},\\
\text{incidence icosaédrique}
&\longrightarrow V_{12}\xrightarrow{P_5}V_5
\xrightarrow{\Gamma_5}\operatorname{STF}_3
\xrightarrow{\Lambda(n)}\mathcal H_{\rm TT}(n),\\
\text{route holonomique enrichie}
&\dashrightarrow G_{\rm hol},\\
\text{carte métrologique}
&\longrightarrow 6.67430\times10^{-11}\,
\mathrm{m^3\,kg^{-1}\,s^{-2}}.
\end{aligned}
\end{equation}

La première ligne exprime une condition intrinsèque et exacte de
périodicité. La deuxième construit
le support et sa réduction transverse sans trace. La troisième conserve une
frontière de construction explicite. La quatrième est une évaluation dans une
carte décimale déclarée : elle identifie la grandeur, mais ne sélectionne pas le caractère
gravitationnel.

\subsection{Droite de grandeur gravitationnelle et échelle chronologique \texorpdfstring{\(108\)}{108}}

Soient \(c_{\rm int}>0\), \(t_0>0\),
\(\ell_0=c_{\rm int}t_0\) et une section positive
\(\hbar_a\) de la droite d'action, avec
\(a\in\{\mathrm{pre},\mathrm{ret}\}\). Le transducteur dimensionnel est

\begin{equation}
\label{rm:eq:gravity-transducer}
G_a(\theta)
=\theta\frac{c_{\rm int}^{5}t_0^2}{\hbar_a}
=\theta\frac{c_{\rm int}^{3}\ell_0^2}{\hbar_a},
\qquad \theta>0.
\end{equation}

L'égalité des deux expressions utilise seulement
\(\ell_0=c_{\rm int}t_0\), et
\([G]=L^3M^{-1}T^{-2}\). La dimension est fixée par la droite ; le
problème mathématique est de sélectionner le caractère adimensionnel \(\theta\).

Le calendrier CT108 définit

\begin{equation}
\label{rm:eq:gravity-clock108}
T_{108}=108t_0,
\qquad
\omega_{108}=\frac{2\pi}{108t_0},
\qquad
R_{108}=\frac{c_{\rm int}}{\omega_{108}}
=\frac{54}{\pi}\ell_0.
\end{equation}

Pour chaque feuillet d'action, soient

\begin{equation}
\label{rm:eq:gravity-clock-energy-mass}
E_{108,a}=\hbar_a\omega_{108},
\qquad
m_{108,a}=\frac{E_{108,a}}{c_{\rm int}^2}.
\end{equation}

Alors, la longueur de Compton réduite coïncide avec le rayon du cycle :

\begin{equation}
\label{rm:eq:gravity-compton-clock}
\bar\lambda_{C,a}
=\frac{\hbar_a}{m_{108,a}c_{\rm int}}
=R_{108}.
\end{equation}

\begin{theorem}[Unicité du sélecteur sous la condition de périodicité HMT]
\label{rm:thm:gravity-circular-selector}
Avec la convention réduite
\(r_g=Gm/c_{\rm int}^2\) et la prémisse de périodicité HMT qui identifie le rayon
gravitationnel du quantum CT108 à son rayon de phase, la condition de clôture
\begin{equation}
\label{rm:eq:gravity-closing-condition}
r_{g,a}(\theta)=R_{108}
\end{equation}
sélectionne une unique solution positive,
\begin{equation}
\label{rm:eq:gravity-theta108}
\boxed{
\theta_{108}^{\rm circ}=\left(\frac{54}{\pi}\right)^2
=295.4525715010569415303525\ldots .}
\end{equation}
Dans cette réalisation,
\begin{equation}
\label{rm:eq:gravity-four-lengths}
R_{108}=\bar\lambda_{C,a}=r_{g,a}=\ell_P^{\rm circ}.
\end{equation}
\end{theorem}

\begin{proof}
Par \cref{rm:eq:gravity-transducer,rm:eq:gravity-clock-energy-mass},
\[
r_{g,a}(\theta)
=\frac{G_a(\theta)m_{108,a}}{c_{\rm int}^2}
=\theta c_{\rm int}t_0^2\omega_{108}
=\theta\frac{\pi}{54}\ell_0.
\]
En le comparant à
\(R_{108}=(54/\pi)\ell_0\), on obtient nécessairement
\(\theta=(54/\pi)^2\). L'équation est linéaire en \(\theta\) et tous les
facteurs sont positifs, donc la solution est unique.
\end{proof}

L'égalité \(r_g=R_{108}\) ne se déduit pas de la dimensionnalité de
\cref{rm:eq:gravity-transducer}. C'est la condition géométrique explicite de cette
réalisation HMT : action, Compton et retour périodique doivent déterminer une même
longueur. Le théorème démontre que, cette condition structurelle étant acceptée, le
caractère \(\theta\) est fixé sans utiliser la valeur SI de \(G\).

\begin{remark}[Convention de rayon]
Le théorème utilise le rayon gravitationnel réduit \(Gm/c^2\). Si on le remplace
par le rayon de Schwarzschild \(2Gm/c^2\), le sélecteur change d'un facteur
deux. Les deux expressions correspondent à des conditions géométriques distinctes,
c'est pourquoi la convention doit être déclarée avant le calcul.
\end{remark}

\subsection{Système de Planck et covariance bidirectionnelle}

\begin{corollary}[Famille de Planck associée à la structure chronologique CT108]
\label{rm:cor:gravity-planck-family}
Pour \(G_a=G_a(\theta_{108}^{\rm circ})\),
\begin{align}
\ell_P&=\frac{54}{\pi}\ell_0,
&t_P&=\frac{54}{\pi}t_0,
\label{rm:eq:gravity-planck-length-time}\\
E_{P,a}&=\frac{\hbar_a}{t_P}
=\frac{\pi\hbar_a}{54t_0}
=\frac{h_a}{108t_0},
&\ell_P^2&=\frac{G_a\hbar_a}{c_{\rm int}^3}
=\theta_{108}^{\rm circ}\ell_0^2.
\label{rm:eq:gravity-planck-energy-area}
\end{align}
Adem\'as,
\begin{equation}
\label{rm:eq:gravity-planck-force-power-density}
F_{P,a}=\frac{c_{\rm int}^4}{G_a},
\qquad
P_{P,a}=\frac{c_{\rm int}^5}{G_a},
\qquad
\rho_{P,a}=\frac{\hbar_a}
{(\theta_{108}^{\rm circ})^2c_{\rm int}^5t_0^4}.
\end{equation}
\end{corollary}

\begin{proof}
Les trois premières identités découlent de
\(\ell_P=R_{108}\), \(t_P=\ell_P/c_{\rm int}\) et
\(h_a=2\pi\hbar_a\). Pour l'aire,
\[
\frac{G_a\hbar_a}{c_{\rm int}^3}
=\theta_{108}^{\rm circ}c_{\rm int}^2t_0^2
=\theta_{108}^{\rm circ}\ell_0^2.
\]
Les expressions restantes sont les compositions dimensionnelles de force,
de puissance et de masse par volume avec la même section.
\end{proof}

Soit
\begin{equation}
\label{rm:eq:gravity-Ract}
R_{\rm act}:=\frac{\hbar_{\rm pre}}{\hbar_{\rm ret}}.
\end{equation}
Comme l'action apparaît au numérateur de la masse chronologique et au
dénominateur de \(G\),
\begin{equation}
\label{rm:eq:gravity-twoway-covariance}
\frac{m_{108,\rm pre}}{m_{108,\rm ret}}=R_{\rm act},
\qquad
\frac{G_{\rm pre}}{G_{\rm ret}}=R_{\rm act}^{-1},
\qquad
G_a\hbar_a
=\theta_{108}^{\rm circ}c_{\rm int}^5t_0^2.
\end{equation}
Par conséquent, \(\ell_P^2=G_a\hbar_a/c^3\) ne change pas de feuillet. La gravité et
l'action s'orientent réciproquement, tandis que l'aire conserve la
géométrie commune. Cette annulation est l'interface de confluence vers l'opérateur
d'aire du \cref{rm:ch:modular}.

\subsection{Composante icosaédrique de dimension \(5\)}

L’antécédent n’est ni \(\R\) ni un choix ultérieur de cinq polarisations. L’objet chargé de transporter l’incidence est défini explicitement avant l’introduction de la réalisation exceptionnelle.

Soit \(X_{\rm TPK}^{\rm exc}\) le sous-espace de l’état \(\TPK\) enrichi
qui conserve, outre la sortie arithmétique, l’orientation, la route,
la retenue et l’incidence exceptionnelle. Pour expliciter l’espace géométrique d’arrivée,
soit \(\varphi_g=(1+\sqrt5)/2\), \(\nu=\sqrt{1+\varphi_g^2}\), et définissons
\begin{equation}
\label{rm:eq:gravity-icosahedron-vertices}
\Omega_{12}^{\rm ico}
=\frac1\nu\bigl(
\{0\}\times\{\pm1\}\times\{\pm\varphi_g\}
\;\cup\;
\{\pm1\}\times\{\pm\varphi_g\}\times\{0\}
\;\cup\;
\{\pm\varphi_g\}\times\{0\}\times\{\pm1\}
\bigr)\subset S^2.
\end{equation}
Ses douze éléments sont les classes orientées de sommet d’un icosaèdre.
L’incidence ne reste pas verbale : elle est fixée par
\begin{equation}
\label{rm:eq:gravity-icosahedron-incidence}
\mathcal I_{12}^{\rm ico}
=\{(v,w)\in\Omega_{12}^{\rm ico}\times\Omega_{12}^{\rm ico}:
v\neq w,\ \langle v,w\rangle=1/\sqrt5\}.
\end{equation}
Chaque sommet a cinq voisins dans cette relation. Le pont de lecture
est typé comme une application
\begin{equation}
\label{rm:eq:gravity-exceptional-reader-map}
\pi_{\rm ico}:X_{\rm TPK}^{\rm exc}\longrightarrow\Omega_{12}^{\rm ico}.
\end{equation}
La cardinalité douze de l’image de \(\pi_{\rm ico}\) ne suffit pas. Pour définir une réalisation d’incidence, l’application doit satisfaire simultanément :
\begin{enumerate}[label=(\roman*)]
\item être surjective et constante exactement sur les fibres qui écartent
la carte arithmétique mais conservent la route exceptionnelle ;
\item entrelacer l’action des automorphismes orientés de l’antécédent
avec une action transitive sur \(\Omega_{12}^{\rm ico}\) ;
\item transporter l’incidence TPK déclarée vers
\(\mathcal I_{12}^{\rm ico}\), et non
ses seuls cardinaux ;
\item conserver l’involution de feuillet et le retour nonadique dans le
stabilisateur de la fibre correspondante.
\end{enumerate}

Ces conditions séparent trois assertions : l’existence de l’application \(\pi_{\rm ico}\), l’identification du groupe qui agit et la représentation linéaire construite ensuite. La deuxième et la troisième ne peuvent servir à définir rétrospectivement la première.

\begin{theorem}[Descente de l’incidence au quotient icosaédrique]
\label{rm:thm:gravity-exceptional-descent}
Supposons que \(\pi_{\rm ico}\) satisfait (i)--(iv) et que l’action
orientée induite sur \(\Omega_{12}^{\rm ico}\) a pour groupe
\(G_{\rm ico}\simeq A_5\). Alors, pour toute classe
\(v_0\in\Omega_{12}^{\rm ico}\),
\begin{equation}
\label{rm:eq:gravity-icosahedral-quotient}
H_{v_0}=\operatorname{Stab}_{G_{\rm ico}}(v_0)\simeq C_5,
\qquad
\Omega_{12}^{\rm ico}\simeq G_{\rm ico}/H_{v_0}\simeq A_5/C_5.
\end{equation}
L’application
\begin{equation}
\label{rm:eq:gravity-coset-map}
\kappa_{v_0}:G_{\rm ico}/H_{v_0}\longrightarrow\Omega_{12}^{\rm ico},
\qquad gH_{v_0}\longmapsto g\cdot v_0,
\end{equation}
est une bijection équivariante et préserve l’incidence si l’on déclare sur le quotient
\begin{equation}
\label{rm:eq:gravity-coset-incidence}
gH_{v_0}\sim g'H_{v_0}
\quad\Longleftrightarrow\quad
(g\cdot v_0,g'\cdot v_0)\in\mathcal I_{12}^{\rm ico}.
\end{equation}
Par conséquent, la composée \(\kappa_{v_0}^{-1}\pi_{\rm ico}\) est l’application
typée de l’antécédent TPK vers le quotient \(A_5/C_5\).
\end{theorem}

\begin{proof}
La liste de \cref{rm:eq:gravity-icosahedron-vertices} a douze éléments.
Par transitivité, \(|G_{\rm ico}\cdot v_0|=12\), et le théorème
orbite--stabilisateur donne \(|H_{v_0}|=60/12=5\). Tout groupe d’ordre
premier cinq est cyclique ; donc \(H_{v_0}\simeq C_5\). La formule de
\(\kappa_{v_0}\) ne dépend pas du représentant : si
\(gH_{v_0}=g'H_{v_0}\), alors \(g^{-1}g'\in H_{v_0}\) et
\(g'v_0=gv_0\). La surjectivité découle de la transitivité, et
l’égalité des cardinaux donne l’injectivité. L’équivariance est immédiate :
\(\kappa_{v_0}(agH_{v_0})=a\kappa_{v_0}(gH_{v_0})\). Enfin,
\cref{rm:eq:gravity-coset-incidence} est précisément le transport de
structure par cette bijection ; l’hypothèse (iii) fait commuter ce
transport avec \(\pi_{\rm ico}\).
\end{proof}

Le théorème ne remplace pas la construction de \(\pi_{\rm ico}\) : il démontre ce qui est obtenu une fois ses quatre conditions vérifiées. Dans le profil HMT complet, le corridor d’incidence qui mène à Paley--Witt--Golay--Leech--\(V^\natural\)--Monstre--Moonshine est le domaine de définition de cette application. Dans ce volume autonome, l’entrée exacte est l’application \cref{rm:eq:gravity-exceptional-reader-map} assortie de ses conditions ; toute la chaîne depuis \(A_5/C_5\) jusqu’à la réduction TT est prouvée ci-dessous sans consulter à nouveau une valeur gravitationnelle.

La chaîne de provenance s’exprime au moyen de tous ses morphismes :
\begin{equation}
\label{rm:eq:gravity-exceptional-provenance}
\APP\longrightarrow\TRIT\longrightarrow\TPK
\longrightarrow X_{\rm TPK}^{\rm exc}
\xrightarrow{\ \pi_{\rm ico}\ }\Omega_{12}^{\rm ico}
\xleftarrow[\simeq]{\ \kappa_{v_0}\ }A_5/C_5.
\end{equation}
Le corridor exceptionnel complet est démontré dans son développement propre. L’hypothèse que ce chapitre ne peut remplacer par un dénombrement est l’existence de \(\pi_{\rm ico}\) satisfaisant (i)--(iv). Cette entrée étant fixée, la descente au quotient est \cref{rm:thm:gravity-exceptional-descent}, et l’on construit à partir d’elle \(P_5\), \(\Gamma_5\) et \(\Lambda(n)\). La réalisation du support résultant comme champ gravitationnel demeure une interface physique typée, et non une conséquence de cardinalité.

Soit \(G=A_5\), soit \(H\simeq C_5\) le stabilisateur d’un sommet
orienté de l’icosaèdre et considérons
\begin{equation}
\label{rm:eq:gravity-v12}
V_{12}=\R[G/H],
\qquad |G/H|=\frac{60}{5}=12.
\end{equation}
La représentation de permutation \(\rho_{12}\) se décompose comme
\begin{equation}
\label{rm:eq:gravity-v12-decomposition}
V_{12}\simeq\mathbf1\oplus V_3\oplus V_{3'}\oplus V_5.
\end{equation}
L’idempotent central du facteur de caractère
\(\chi_5=(5,1,-1,0,0)\) est
\begin{equation}
\label{rm:eq:gravity-p5}
\boxed{
P_5=\frac5{60}\sum_{g\in A_5}\chi_5(g^{-1})\rho_{12}(g)
=\frac1{12}\left(
5I+\sum_{g\in2A}\rho_{12}(g)
-\sum_{g\in3A}\rho_{12}(g)
\right).}
\end{equation}

\begin{theorem}[Projecteur et couplage gravitationnel de rang cinq]
\label{rm:thm:gravity-p5-coupling}
L'opérateur \(P_5\) satisfait
\begin{equation}
\label{rm:eq:gravity-p5-properties}
P_5^2=P_5=P_5^*,
\qquad \Tr P_5=\rank P_5=5.
\end{equation}
Pour
\begin{equation}
\label{rm:eq:gravity-G5}
\mathbb G_{5,a}:=G_a(\theta_{108}^{\rm circ})P_5
\end{equation}
on a
\begin{equation}
\label{rm:eq:gravity-G5-spectrum}
\Spec(\mathbb G_{5,a})
=\{G_a(\theta_{108}^{\rm circ})\text{ avec multiplicité }5,
\;0\text{ avec multiplicité }7\}.
\end{equation}
\end{theorem}

\begin{proof}
L’orthogonalité des caractères fait de \cref{rm:eq:gravity-p5}
l’idempotent central qui agit comme l’identité sur l’unique copie de
\(V_5\) et comme zéro sur les trois autres facteurs de
\cref{rm:eq:gravity-v12-decomposition}. La réalité du caractère et
l’orthogonalité de \(\rho_{12}\) donnent l’autoadjonction. Sa trace et son rang
sont cinq. Multiplier un projecteur de rang cinq par le scalaire positif
\(G_a\) produit immédiatement le spectre indiqué.
\end{proof}

L’expression \(\mathbb G_{5,a}\) n’ajoute pas cinq valeurs différentes de la constante gravitationnelle. C’est une même section de \(G\) agissant sur le facteur tensoriel sélectionné et annulant les sept autres modes du module de permutation.


\paragraph{Réduction équivariante du support au secteur tensoriel sans trace.} Le support est entrelacé avec les tenseurs symétriques sans trace et réduit de manière équivariante avant d’imposer la transversalité.
